\section{What This Report Is Trying To Do}
This report explains the memory estimator that is currently implemented in
SimpleSFT. It is not a historical survey and it is not a proposal for a future
model. The goal is practical:

\begin{quote}
Given a model and a training configuration, estimate the peak GPU memory used
by one rank during one training step.
\end{quote}

The estimator is easiest to understand if you stop thinking about ``VRAM'' as
one giant blob. The code does not predict one opaque number. It builds memory
from a small set of named objects, then asks which training phase is worst.

\section{The One-Step Mental Model}
For memory purposes, one training step has three moments that matter:
\begin{enumerate}[leftmargin=1.5em]
\item \textbf{Forward:} parameters are resident, some activations are kept for
backward, and kernels use temporary workspace.
\item \textbf{Backward:} gradients now exist, some forward state is still
alive, and backward kernels need their own workspace.
\item \textbf{Optimizer step:} gradients and optimizer state remain live while
the optimizer fetches parameters, updates them, and communicates with other
ranks if training is distributed.
\end{enumerate}

The estimator assigns one candidate peak to each of those moments and then
takes the maximum:
\[
M_{\mathrm{peak}} =
\max\!\left(
M_{\mathrm{forward}},
M_{\mathrm{backward}},
M_{\mathrm{optimizer}}
\right).
\]

That is the whole top-level idea. Everything else in this report explains where
those three candidate peaks come from.

\section{The Four Memory Families}
Almost every estimator choice fits into one of four buckets:
\begin{itemize}[leftmargin=1.5em]
\item \textbf{Resident state:} memory that simply exists on the rank, such as
parameters, gradients, optimizer state, master weights, and fixed backend
support buffers.
\item \textbf{Retained activations:} forward-pass tensors that are intentionally
kept alive because backward still needs them.
\item \textbf{Workspace:} temporary memory that spikes while a phase is
actively running.
\item \textbf{Carryover:} state that is not part of the clean resident floor
but is still hanging around when the optimizer step begins.
\end{itemize}

This decomposition matters because different training regimes move memory
between these buckets. Checkpointing, for example, reduces retained activation
state but usually increases backward-time recompute workspace. ZeRO changes the
resident state and the optimizer-step windows without changing the model's
architecture. LoRA changes which parameters and linears are actually trainable.

\section{Why The Estimator Works Per Rank}
The estimator always answers a per-rank question. That is the right unit for
real training failures because OOM happens on a device, not on the cluster in
the abstract.

Two consequences follow immediately:
\begin{itemize}[leftmargin=1.5em]
\item Tensor parallelism changes the local width of many tensors.
\item Data parallel sharding changes how much of a resident or optimizer object
one rank owns.
\end{itemize}

So two jobs with the same model and sequence length can have very different
per-rank peaks if they use different world sizes, tensor-parallel degrees, or
ZeRO stages.

\section{What The Estimator Takes As Input}
The estimator combines two inputs:
\begin{itemize}[leftmargin=1.5em]
\item an \texttt{EstimatorConfig}, which describes how training is being run,
\item a \texttt{ModelSpec}, which describes the model's size and inspected
layer structure.
\end{itemize}

\paragraph{The configuration tells us.}
\begin{itemize}[leftmargin=1.5em]
\item whether we are doing full fine-tuning or LoRA,
\item whether training is single-GPU, DDP, ZeRO-2, or ZeRO-3,
\item which optimizer family is used,
\item which attention backend is active,
\item whether checkpointing is enabled,
\item the world-size, tensor-parallel, and sequence-parallel choices,
\item micro-batch size, sequence length, and dtypes,
\item explicit runtime support assumptions and distributed bucket sizes.
\end{itemize}

\paragraph{The model spec tells us.}
\begin{itemize}[leftmargin=1.5em]
\item the total parameterization,
\item hidden, intermediate, and vocabulary sizes,
\item attention widths and head counts,
\item which linears are trainable,
\item and, when available, explicit parameter metadata for more exact tensor
accounting.
\end{itemize}

\section{A Small Notation Guide}
Only a few symbols matter for the rest of the report:
\begin{itemize}[leftmargin=1.5em]
\item $L$: number of transformer layers.
\item $H$: hidden size.
\item $I$: intermediate size.
\item $V$: vocabulary size.
\item $B$: micro-batch size per rank.
\item $S$: sequence length.
\item $T = B \cdot S$: tokens processed by one rank before any sequence
parallel split.
\item $T_{\mathrm{TP}}$: tensor-parallel degree.
\item $D_{\mathrm{P}}$: data-parallel degree.
\end{itemize}

When the report refers to a byte width such as $b_w$ or $b_g$, it simply means
``bytes per element for that dtype.'' For example, bf16 weights imply
$b_w = 2$.

\section{How To Read The Rest}
The remainder of the report follows the same order as the estimator itself:
\begin{enumerate}[leftmargin=1.5em]
\item build the resident memory floor,
\item add the retained activations that survive the forward pass,
\item add the phase-local workspace,
\item assemble the forward, backward, and optimizer-step candidates,
\item pick the largest one.
\end{enumerate}

Whenever a term is derived directly from concrete tensors, the report says so.
Whenever a term is really a runtime or backend assumption, the report says that
too. That distinction is important. It keeps the model honest.
